\newpage

\section{\huge \textcolor{blue}{傅里叶级数}}
\subsection{\LARGE 共振}
\subsubsection{\large 无阻尼}
下文利用共振引入相应学习点, 更着重于数学形式, 而非物理概念\\
下先给出输入为$\cos \omega_1 t$方程:
\[ y'' + \omega_0^2 y = \cos \omega_1 t \]
\[ (D^2 + \omega_0^2) y = \cos \omega_1 t \]
复化方程:
\[ (D^2 + \omega_0^2) \tilde{y} = e^{i \omega_1 t} \]
\[\tilde{y} = \frac{e^{i \omega_1 t}}{(i \omega_1)^2 + \omega_0^2} = \frac{e^{i \omega_1 t}}{\omega_0^2 - \omega_1^2}\]
所以响应$y_p$为:
\[ y_p = Re(\tilde{y}) = \frac{\cos \omega_1 t}{\omega_0^2 - \omega_1^2}\]
得:\\
当$\omega_1  \to  \omega_0$时, $y_p \to \infty $, 即
输入频率越接近物体的固有频率时, 振幅越大.
\begin{figure}[htbp]
    \centering
\includegraphics[width=0.5\textwidth]{C:/Users/37/Desktop/Notes on Common Calculus Equations/Notes and pictures/cwf9.1.1.png}
    \caption{$\omega_1  \to  \omega_0$}
    \label{fig:cwf9.1.1}
\end{figure}

若$\omega_1 = \omega_0$
\[ (D^2 + \omega_0^2) y = \cos \omega_0 t \]
\[ (D^2 + \omega_0^2) \tilde{y} = e^{i \omega_0 t} \]
所以 $i \omega_0$是$D^2 + \omega_0^2$的一个单根
\[ \tilde{y} = \frac{t e^{i \omega_0 t}}{2 i \omega_0} \]
\[ y_p = Re(\tilde{y}) = \frac{t \sin \omega_0 t}{2 \omega_0} \]
\begin{figure}[htbp]
    \centering
\includegraphics[width=0.6\textwidth]{C:/Users/37/Desktop/Notes on Common Calculus Equations/Notes and pictures/cwf9.1.2.png}
    \caption{$\omega_1 = \omega_0$}
    \label{fig:cwf9.1.2}
\end{figure}

在 $\omega_1 \to \omega_0$时, $0$处的振幅值逐渐增大, 
在 $\omega_1 = \omega_0$时, 值为$0$, 这之间的变化如何解释?\\
不是唯一特解, 要知任何满足非齐次方程的解都是一个特解
\[ (D^2 + \omega_0^2) y = \cos \omega_1 t \]
在特解上加上通解:
\[ y =  \frac{\cos \omega_1 t}{\omega_0^2 - \omega_1^2} - \frac{\cos \omega_0 t}{\omega_0^2 - \omega_1^2} \] 
现对此形式求极限:
\begin{align*}
    \lim_{\omega_1 \to \omega_0} \frac{\cos \omega_1 t - \cos \omega_0 t}{\omega_0^2 - \omega_1^2} &= \lim_{\omega_1 \to \omega_0} \frac {-t \sin \omega_1 t}{- 2 \omega_1} \\
    &= \frac{t \sin \omega_0 t}{2 \omega_0}      
\end{align*}

得另一形式解, 继续探求联系:\\
知:
\[ \cos B - \cos A = 2 \sin \frac{A - B}{2} \sin \frac{A + B}{2} \]
有:
\[ \frac{\cos \omega_1 t - \cos \omega_0 t}{\omega_0^2 - \omega_1^2} = \frac{2 \sin \frac{\omega_0 - \omega_1}{2} t \sin \frac{\omega_1 + \omega_0}{2} t}{\omega_0^2 - \omega_1^2} \]
将$\frac{2 \sin \frac{\omega_0 - \omega_1}{2} t}{\omega_0^2 - \omega_1^2}$视为振幅, $sin \frac{\omega_1 + \omega_0}{2} t$视为基本振荡
\begin{figure}[htbp]
    \centering
\includegraphics[width=0.6\textwidth]{C:/Users/37/Desktop/Notes on Common Calculus Equations/Notes and pictures/cwf9.1.3.png}
    \caption{联系图}
    \label{fig:cwf9.1.3}
\end{figure}

如图:展示了频率变化的曲线关系, 而随$\omega_1 \to \omega_0$ , 振幅值越来越大, 越不过图像最大值,使得图像显示为此图前半部分.

\subsubsection{\large 阻尼}
仅先作了解, 
有方程:
\[ m x'' + c x' + k x = f(t) \]
\[ x'' + 2 p x' + \omega_0^2 x = f(t) \]
有频率关系:
\[ \omega_1^2 = p^2 + \omega_0^2 \]
若在输入引入新频率:
\[ y'' + 2 p y' + \omega_0^2 y = \cos \omega t \]
思考:什么样的输入频率, 才能使响应的振幅最大?
\[ \omega_r = \sqrt{\omega_0^2 - 2p^2} \]

\subsection{\LARGE 傅里叶级数引入}
在方程:
\[ y''  + a y' + b y = f(t) \]
考虑输入$f(t)$为 $e^{a t}, \sin t, \cos t, \cdots \cdots  $\\
且任何平常的周期函数, $T = 2 \pi$, 可表示为:
\[ f(t) = c_0 + \sum_{n = 1}^{\infty} (a_n \cos nt  + b_n \sin nt) \] 
\begin{align*}
   \text{输入}  &\to \text{响应}  \\
   b_n \sin nt  &\to b_n y_n^{s}(t)   \\
   a_n \cos nt  &\to a_n y_n^{c}(t) \\      
\end{align*}
两侧求和得:
\[ f(t) \to   c + \sum_{n = 1}^{\infty} (a_n y_n^{c}(t) + b_n y_n^{s}(t)) \]
能如此附带系数相加, 因有叠加原理, 更一般的原因为ODE是线性的.\\
下继续完善傅里叶级数:\\
给定周期函数$f(t)$, $  T = 2 \pi $\\
先考虑正交性, 有下:\\
若$u(t), v(t), T = 2 \pi, t \in \mathbb{R} $, 
$u , v$在 $[- \pi ,  \pi]$上正交, 则:
\[ \int_{- \pi}^{\pi} u(t) v(t) \,dt = 0 \] 
同时有如下性质:
\[ \left\{ \sin nt , \cos mt \right\} , n = 1, 2, 3, \cdots \cdots  , m = 0, 1, 2, \cdots \cdots   \] 
此集合中任意两不同元素均在区间上正交, 除下:
\[
\int_{- \pi}^{\pi} \sin^2 nt \,dt = \int_{- \pi}^{\pi} \cos^2 mt \,dt = \pi
\]
关于此性质证明有大概如下3种方法:
\begin{enumerate}
    \item 三角恒等式变形
    \item 化为复指数
    \item ODE本身(推广性强)
\end{enumerate}
对此只展示法3:\\
证明:\\
假设$m \neq n$, 满足的微分方程为:
\[ u'' + n^2 u = 0 \]
取$u_n, v_m$为任意的两个函数, 变形:
\[ u'' = -n^2 u \]
作:
\begin{align*}
   \int_{- \pi}^{\pi} u_n'' v_m \,dt &= u_n' v_m \vert _{-\pi}^{\pi} - \int_{- \pi}^{\pi} u_n' v_m' \,dt  \\
    &= - \int_{- \pi}^{\pi} u_n' v_m' \,dt   \\
    &= - n^2 \int_{- \pi}^{\pi}  u_n v_m \,dt     
\end{align*}
其中$u_n' v_m \vert _{-\pi}^{\pi} = 0$ , 证明较为繁琐, 从处先使用, 后续可查阅相应资料, 并同理:
\[ \int_{- \pi}^{\pi} v_m'' u_n \,dt = - m^2 \int_{- \pi}^{\pi}  v_m u_n \,dt \]
从中观察, 知$u_n',  v_m'$对称, $u_n, v_m$非对称, 且$m \neq n$, 故:\\
$u_n, v_m$只能在$[- \pi, \pi]$上正交, 才能使等式成立.\\
下求傅里叶级数系数:\\
给定$f(t), t \in [- \pi, \pi] , T = 2 \pi $

\begin{align*}
                                     f(t) &= \cdots + a_k \cos kt +  \cdots  + a_n \cos nt +  \cdots  \\
    \int_{- \pi}^{\pi}  f(t) \cos nt \,dt &= \cdots   + \int_{- \pi}^{\pi} a_k \cos kt  \cos nt \,dt + \cdots + \int_{- \pi}^{\pi} a_n \cos^2 nt  \,dt   \\
                                          &= a_n \pi
\end{align*}
解得:
\begin{align*}
   a_n  &= \frac{1}{\pi} \int_{- \pi}^{\pi}  f(t) \cos nt \,dt \hspace*{1em} , n = 1, 2, 3, \cdots  \\
   b_n  &= \frac{1}{\pi} \int_{- \pi}^{\pi}  f(t) \sin nt \,dt \hspace*{1em} , n = 1, 2, 3, \cdots     
\end{align*}
求常数项:
\begin{align*}
                            f(t)  &= c_0 + \cdots  +  a_n \cos nt +  \cdots  \\
    \int_{- \pi}^{\pi}  f(t) \,dt &= 2\pi c_0  + \cdots + \int_{- \pi}^{\pi} a_n \cos nt  \,dt   \\
                             c_0  &= \frac{1}{2\pi} \int_{- \pi}^{\pi}  f(t) \,dt
\end{align*}
为了让公式统一, 令$c_0 = a_0$, 即:
\[ a_n = \frac{1}{\pi} \int_{- \pi}^{\pi}  f(t)  \,dt \]


\subsection{\LARGE 傅里叶级数拓展}
整理前文:
\[ f(t) = \frac{a_0}{2} + \sum_{n = 1}^{\infty} (a_n \cos nt  + b_n \sin nt) \]
\begin{align*}
   a_n  &= \frac{1}{\pi} \int_{- \pi}^{\pi}  f(t) \cos nt \,dt \hspace*{1em} , n = 1, 2, 3, \cdots  \\
   b_n  &= \frac{1}{\pi} \int_{- \pi}^{\pi}  f(t) \sin nt \,dt \hspace*{1em} , n = 1, 2, 3, \cdots     
\end{align*}
有:给定$a_n, b_n$ , $f(t)$确定,即\\
一个函数$f(t)$不能有两个不同的傅里叶级数\\
下针对傅里叶级数的研究有些拓展方向, 如下:
\subsubsection{\large 精简计算}
此方向针对奇偶函数:\\
若$f(t)$为偶函数, 有$f(-t) = f(t) $, 则:\\
$f(t)$的傅里叶级数只包括 cos 项
\[ f(t) = \frac{a_0}{2} + \sum_{n = 1}^{\infty} a_n \cos nt  , \hspace*{1em} (b_n = 0)\]
证明:
\[ f(t) = \frac{a_0}{2} + \sum_{n = 1}^{\infty} (a_n \cos nt  + b_n \sin nt)  \]
\[ f(-t) = \frac{a_0}{2} + \sum_{n = 1}^{\infty} (a_n \cos nt  - b_n \sin nt) \]
$ \because \hspace*{1em} f(t) = f(-t) $\\
通过对比系数
\[ \therefore \hspace*{1em} b_n = 0 \]
同, 若$f(t)$为奇函数, 则 $a_n = 0$, 即
\[ f(t) =  \sum_{n = 1}^{\infty} b_n \sin nt  \]
进一步对积分化简:\\
若$f(t)$为偶函数, 则 $f(t)\cos nt $ 同为偶函数
\[ a_n  = \frac{2}{\pi} \int_{0}^{\pi}  f(t) \cos nt \,dt ,\hspace*{1em} (b_n = 0)  \]
若$f(t)$为奇函数, 则 $f(t)\sin nt $ 为偶函数
\[ b_n  = \frac{2}{\pi} \int_{0}^{\pi}  f(t) \sin nt \,dt  \]

\subsubsection{\large 补充}
在分析中了解:\\
泰勒级数是用来分析一个函数靠近某点是的表现, \\
傅里叶级数是用来分析一个函数在某区间内的表现.\\
下再介绍一个重要定理:\\
若$f$是一个连续函数, 且在$t_0$点附近连续, 则:\\
傅里叶级数是收敛的, 且
$f(t)$的值为它的傅里叶级数在$t_0$处的和;\\
若$t_0$是跳跃性不连续点, 那么\\
傅里叶级数相加是收敛的, 且为跳跃处的中点

\subsubsection{\large 拓展周期}
现设: $T = 2 L$, 对坐标轴伸缩变化:
\[ f(t) = \frac{a_0}{2} + \sum_{n = 1}^{\infty} (a_n \cos \frac{n \pi}{L}t  + b_n \sin \frac{n \pi}{L}t )  \]
并整理上述结论并运用:\\
若$f(t)$为偶函数, 则:
\[ a_n  = \frac{2}{\pi} \int_{0}^{2L}  f(t) \cos \frac{n \pi}{L}t \,dt ,\hspace*{1em} (b_n = 0)  \]

\subsubsection{\large 拓展非周期函数}
如题, 对于有限区间的函数均可使用傅里叶级数表达\\
需对本体区间扩展, 将整个区间嵌入一个周期内\\
eg: $y(t) = t^2$, $t \in [0, L] $
\begin{figure}[htbp]
    \centering
\includegraphics[width=0.4\textwidth]{C:/Users/37/Desktop/Notes on Common Calculus Equations/Notes and pictures/cwf9.2.1.png}
    \caption{原区间图}
    \label{fig:cwf9.2.1}
\end{figure}

现在对此进行奇偶拓展, 如图:
\begin{figure}[htbp]
    \centering
\includegraphics[width=0.5\textwidth]{C:/Users/37/Desktop/Notes on Common Calculus Equations/Notes and pictures/cwf9.2.2.png}
    \caption{奇偶拓展图}
    \label{fig:cwf9.2.2}
\end{figure}

如上选取一种拓展后, 按正常步骤处理即可.
\subsection{\LARGE 利用傅里叶级数求特解}
设一方波: $f(t)$ , 非奇非偶, $L = 1$, 如下:
\begin{figure}[htbp]
    \centering
\includegraphics[width=0.5\textwidth]{C:/Users/37/Desktop/Notes on Common Calculus Equations/Notes and pictures/cwf9.2.3.png}
    \caption{方波图}
    \label{fig:cwf9.2.3}
\end{figure}

为更方便求解, 进行必要变形, 将其向下偏移$\frac{1}{2}$, 
此时为奇函数, 记为 $s(t)$:

\begin{figure}[htbp]
    \centering
\includegraphics[width=0.5\textwidth]{C:/Users/37/Desktop/Notes on Common Calculus Equations/Notes and pictures/cwf9.2.4.png}
    \caption{平移图}
    \label{fig:cwf9.2.4}
\end{figure}

为数值方便计算, 将轴伸缩变化 $u = \pi t$
\begin{figure}[htbp]
    \centering
\includegraphics[width=0.5\textwidth]{C:/Users/37/Desktop/Notes on Common Calculus Equations/Notes and pictures/cwf9.2.5.png}
    \caption{轴伸缩图}
    \label{fig:cwf9.2.5}
\end{figure}

计算如下:
\[ g(u) = \frac{4}{\pi} \sum_{n = odd}^{\infty}\frac{\sin n u}{n} \]
\[ f(t) = s(t) + \frac{1}{2} \]
\[ s(t) = \frac{1}{2} g(u) \]
\[ f(t) = \frac{1}{2} g(u) + \frac{1}{2} \]
\[ \therefore  \hspace*{1em} f(t) = \frac{1}{2} + \frac{2}{\pi} \sum_{n = odd}^{\infty}\frac{\sin n \pi t}{n} \]
到现在, 为了解决二阶线性常系数方程
\[ x'' + \omega^2 x = f(t) \]
求$x_p$\\
若$f(t) = \left\{
                \begin{aligned}
                \sin \omega t  \\
                \cos \omega t
                \end{aligned}
\right.$ , 解为:
\[x_p =\frac{\left\{
                \begin{aligned}
                 \sin \omega x  \\
                 \cos \omega x
                \end{aligned}
                \right.}{\omega_0^2 - \omega^2} 
\]
即:
\[ f(t) = \frac{a_0}{2} + \sum_{n = 1}^{\infty} (a_n \cos \omega_n t  + b_n \sin \omega_n t )  \]
其中$\omega_n = \frac{n \pi}{L}, T = 2L$\\
解:
\[
x_p = \frac{a_0}{2 \omega_0^2} + \sum_{n = 1}^{\infty} ( \frac{a_n \cos \omega_n t}{\omega_0^2 - \omega^2} + \frac{b_n \sin \omega_n t}{\omega_0^2 - \omega^2})
\]
若输入$f(t)$为上文中的方波, 解得

\[
x_p = \frac{1}{2 \omega_0^2} + \frac{2}{\pi} \sum_{n = odd}^{\infty} \frac{\sin n \pi t}{n (\omega_0^2 - (n \pi)^2)}
\]

下再介绍一代换法: \\
将傅里叶级数逐项微分, 得系数\\
假设$f(t)$ 解的形式为:
\[ x_p = c_0 + \sum_{n = 1}^{\infty} c_n \sin n \pi t  \]
\[ x_p'' = - \sum_{n = 1}^{\infty} c_n (n \pi)^2 \sin n \pi t  \]
代入ODE, 
\[ x_p'' + \omega^2 x_p = f(t) \]
通过系数等价, 即求得$c_n$.